\chapter{Cài đặt, kiểm thử và triển khai hệ thống}
\label{chap:experiment}

\section{Môi trường cài đặt và triển khai}

\subsection{Môi trường phát triển}

Fitty được phát triển dưới dạng ứng dụng Android bằng Kotlin và Gradle Kotlin DSL. Môi trường phát triển cần hỗ trợ Jetpack Compose, Firebase, Room, Hilt và WorkManager, phù hợp với tập công nghệ đã trình bày ở chương trước \cite{android_jetpack,android_compose,android_hilt,android_room,android_workmanager,firebase_setup}. Các thông số chính được tóm tắt trong Bảng~\ref{tab:dev-environment}.

\begin{table}[H]
\centering
\small
\caption{Thông số môi trường phát triển và triển khai của hệ thống}
\label{tab:dev-environment}
\begin{tabular}{|p{4.3cm}|p{7.8cm}|}
\hline
\textbf{Thành phần} & \textbf{Thông tin} \\ \hline
Nền tảng ứng dụng & Android \\ \hline
Ngôn ngữ lập trình & Kotlin \\ \hline
Kiểu dự án build & Gradle Kotlin DSL \\ \hline
SDK biên dịch & compileSdk 36 \\ \hline
Thiết bị tối thiểu & minSdk 24 \\ \hline
Kiểu giao diện & Jetpack Compose + Material 3 \\ \hline
Lưu trữ cục bộ & Room, DataStore \\ \hline
Dịch vụ phía sau & Firebase Auth, Firestore, Storage, Messaging, Analytics \\ \hline
Tác vụ nền & WorkManager \\ \hline
Media và ảnh & Coil, Media3/ExoPlayer \\ \hline
\end{tabular}
\normalsize
\end{table}

Trong quá trình cài đặt, các tham số cấu hình môi trường và thông tin kết nối dịch vụ được tách riêng khỏi phần hiện thực ứng dụng để thuận tiện cho triển khai và vận hành.

\subsection{Môi trường chạy ứng dụng}

Về vận hành, ứng dụng được thiết kế để chạy trên thiết bị Android có Google Play Services khi cần dùng xác thực Google hoặc Firebase Cloud Messaging. Môi trường chạy cũng cần kết nối đúng Firebase project để bảo đảm các luồng đồng bộ, meal scan, body scan và thông báo hoạt động đúng \cite{firebase_auth,firebase_fcm}.

\section{Định hướng hiện thực}

Phần cài đặt của hệ thống được triển khai bám theo kiến trúc đã trình bày ở chương trước, trong đó giao diện người dùng, xử lý nghiệp vụ, dữ liệu cục bộ và dịch vụ phía sau được tách thành các lớp trách nhiệm tương đối rõ ràng. Cách hiện thực này giúp việc mở rộng chức năng và kiểm thử từng luồng chính thuận lợi hơn, đồng thời tránh để logic nghiệp vụ bị phân tán trực tiếp vào giao diện.

\section{Hiện thực các phân hệ chính}

\subsection{Phân hệ truy cập và khởi động}

Phân hệ truy cập hỗ trợ đăng ký, đăng nhập bằng email, đăng nhập Google, dùng thử ở chế độ khách và khôi phục mật khẩu. Sau bước xác thực, hệ thống đồng bộ trạng thái hồ sơ người dùng và xác định màn hình khởi động phù hợp để điều hướng tới onboarding hoặc vùng chức năng chính \cite{firebase_auth,firebase_firestore}.

\subsection{Phân hệ onboarding và kế hoạch tập khởi đầu}

Luồng onboarding thu thập thông tin nền của người dùng như mục tiêu tập luyện, chiều cao, cân nặng, cân nặng mục tiêu, số ngày tập, thời gian tập, điều kiện dụng cụ và định hướng dinh dưỡng. Từ dữ liệu này, hệ thống hình thành kế hoạch tập khởi đầu, lịch buổi tập và nội dung xem trước kế hoạch. Trong quá trình hiện thực, phần thu thập dữ liệu và phần xử lý hình thành kế hoạch được tách biệt để thuận lợi cho bảo trì và kiểm thử.

\subsection{Phân hệ nội dung bài tập và buổi tập}

Phân hệ bài tập được hiện thực theo hướng ngoại tuyến ưu tiên: dữ liệu mô tả được đồng bộ từ nguồn từ xa, chuẩn hóa và lưu cục bộ trước khi phục vụ giao diện \cite{firebase_firestore,android_room,android_workmanager}. Ở mức nội dung, mỗi bài tập có thể gắn với nhóm cơ, mô tả, bước thực hiện, lỗi sai thường gặp, mẹo thực hiện và media minh họa. Phân hệ buổi tập hỗ trợ cả buổi tập theo lịch và buổi tập nhanh; khi hoàn thành, hệ thống cập nhật kết quả buổi tập, tiến độ và trạng thái lịch tập liên quan.

\subsection{Phân hệ theo dõi tiến độ, dinh dưỡng, hồ sơ và thông báo}

Phân hệ theo dõi tiến độ tổng hợp các thông tin theo ngày, dữ liệu dinh dưỡng, số đo cơ thể, lịch sử body scan và thống kê luyện tập để hiển thị mức độ tiến triển chung. Ở mức hiện thực, hệ thống hỗ trợ ghi nhận nhật ký bữa ăn, phân tích ảnh bữa ăn, lưu dữ liệu body scan và hiển thị một số chỉ số như BMI, cân nặng mục tiêu, lượng calo và thành phần dinh dưỡng cơ bản. Phân hệ hồ sơ hỗ trợ cập nhật thông tin cá nhân và ảnh đại diện qua Firebase Storage. Phân hệ thông báo kết hợp thông báo cục bộ với FCM để duy trì tương tác với người dùng \cite{firebase_storage,firebase_fcm}.

\subsection{Phân hệ AI coach và hỗ trợ phân tích}

Ngoài các luồng nghiệp vụ truyền thống, hệ thống còn hiện thực một số thành phần AI hỗ trợ trải nghiệm người dùng. Cụ thể, ứng dụng có thể gửi câu hỏi tới AI coach để nhận phản hồi ngắn gọn về tập luyện hoặc dinh dưỡng, đồng thời sử dụng mô hình phân tích ảnh cho bài toán nhận diện bữa ăn và đánh giá body scan ở mức hỗ trợ. Phân hệ này không thay thế chuyên gia thể hình hay chuyên gia dinh dưỡng, nhưng đóng vai trò tăng tính tương tác và giúp người dùng phổ thông có thêm gợi ý ban đầu trong quá trình sử dụng.

\section{Triển khai nội dung động}

Một đặc điểm đáng chú ý của hệ thống là nhiều phần nội dung được quản lý từ dịch vụ dữ liệu phía sau thay vì cố định trong ứng dụng, chẳng hạn nội dung onboarding, danh mục bài tập, mẫu kế hoạch khởi đầu và một số cấu hình hành vi. Cách tổ chức này giúp việc cập nhật nội dung linh hoạt hơn, đồng thời giảm nhu cầu phát hành lại ứng dụng khi chỉ thay đổi dữ liệu vận hành. Đối với một ứng dụng hướng tới người tập gym, điều này đặc biệt quan trọng vì nhóm bài tập, gợi ý nội dung và chiến lược trình bày có thể cần điều chỉnh thường xuyên theo mục tiêu vận hành sản phẩm.

\section{Kiểm thử hệ thống}

\subsection{Kiểm thử đơn vị}

Dự án có các bài kiểm thử đơn vị tập trung vào những khu vực có logic nghiệp vụ quan trọng như hình thành kế hoạch tập khởi đầu, đồng bộ dữ liệu bài tập, hoàn thành buổi tập, cập nhật dữ liệu theo ngày, ghi nhận dinh dưỡng và quản lý dữ liệu cục bộ theo từng người dùng. Các bài kiểm thử này giúp xác nhận độ đúng của các khối xử lý trọng yếu trước khi đánh giá ở mức chức năng \cite{android_jetpack,android_room}.

\begin{table}[H]
\centering
\small
\caption{Các nhóm kiểm thử đơn vị tiêu biểu trong dự án}
\label{tab:unit-tests}
\begin{tabular}{|p{4.7cm}|p{7.4cm}|}
\hline
\textbf{Nhóm kiểm thử} & \textbf{Mục tiêu chính} \\ \hline
Sinh kế hoạch khởi đầu & Kiểm tra logic hình thành kế hoạch và nội dung xem trước theo hồ sơ người dùng. \\ \hline
Gợi ý và phân loại bài tập & Kiểm tra quy tắc lựa chọn bài tập theo mục tiêu, mức độ tập và điều kiện thiết bị. \\ \hline
Đồng bộ dữ liệu bài tập & Kiểm tra việc kiểm tra hợp lệ và phân loại trạng thái đồng bộ dữ liệu bài tập. \\ \hline
Hoàn thành buổi tập & Kiểm tra luồng ghi nhận kết quả buổi tập và cập nhật lịch tập liên quan. \\ \hline
Theo dõi tiến độ & Kiểm tra tính toán tiến độ và một số giá trị tổng hợp trên màn hình theo dõi. \\ \hline
Kế hoạch và nội dung động & Kiểm tra việc sử dụng dữ liệu nội dung và nhóm bài tập trong các luồng kế hoạch. \\ \hline
Dữ liệu cục bộ theo người dùng & Kiểm tra phạm vi dữ liệu cục bộ để tránh lẫn giữa các phiên người dùng khác nhau. \\ \hline
Thông báo & Kiểm tra cấu hình và hành vi quản lý thông báo cục bộ. \\ \hline
Dinh dưỡng và tổng hợp theo ngày & Kiểm tra việc ghi nhận dữ liệu dinh dưỡng và cập nhật tổng hợp theo ngày. \\ \hline
\end{tabular}
\normalsize
\end{table}

\subsection{Kiểm thử chức năng}

Ở mức chức năng, kiểm thử tập trung vào các luồng sử dụng chính: truy cập tài khoản, hoàn tất khởi tạo hồ sơ, kích hoạt kế hoạch tập khởi đầu, duyệt nội dung bài tập, thực hiện buổi tập, theo dõi tiến độ và nhận thông báo. Cách tiếp cận này giúp đánh giá được các luồng liên phân hệ thay vì chỉ kiểm tra từng màn hình riêng lẻ.

\begin{table}[H]
\centering
\small
\caption{Các kịch bản kiểm thử chức năng tiêu biểu}
\label{tab:functional-test-scenarios}
\begin{tabular}{|p{1.1cm}|p{3.9cm}|p{4.5cm}|p{2.5cm}|}
\hline
\textbf{Mã} & \textbf{Kịch bản kiểm thử} & \textbf{Kết quả mong đợi} & \textbf{Mức độ} \\ \hline
T1 & Đăng nhập bằng email và mật khẩu hợp lệ & Người dùng vào màn hình chính hoặc khởi tạo hồ sơ tùy trạng thái hiện tại & Cao \\ \hline
T2 & Hoàn tất khởi tạo hồ sơ và mở phần xem trước kế hoạch & Hệ thống lưu hồ sơ, hình thành kế hoạch khởi đầu và hiển thị nội dung xem trước & Cao \\ \hline
T3 & Mở kế hoạch tập và chọn một nhóm bài tập & Danh sách bài tập được lọc đúng theo nhóm cơ & Cao \\ \hline
T4 & Đồng bộ dữ liệu bài tập khi có mạng & Dữ liệu mô tả bài tập được cập nhật vào bộ nhớ cục bộ và trạng thái đồng bộ thành công & Cao \\ \hline
T5 & Mở danh sách bài tập khi mất mạng nhưng đã có bộ nhớ đệm & Danh sách bài tập vẫn hiển thị từ dữ liệu cục bộ & Cao \\ \hline
T6 & Bắt đầu và hoàn thành buổi tập nhanh & Kết quả buổi tập được lưu, tổng hợp theo ngày và tổng số buổi tập được cập nhật & Cao \\ \hline
T7 & Ghi nhận nhật ký bữa ăn hoặc xác nhận kết quả phân tích ảnh bữa ăn & Chỉ số dinh dưỡng, tổng calo và số bữa ăn trong ngày được làm mới & Trung bình \\ \hline
T8 & Cập nhật ảnh đại diện trong hồ sơ & Ảnh được lưu lên dịch vụ lưu trữ, giao diện hiển thị ảnh mới & Trung bình \\ \hline
T9 & Gửi thử thông báo đẩy tới thiết bị đang mở ứng dụng & Banner hoặc thông báo hệ thống xuất hiện và có bản ghi thông báo cục bộ & Cao \\ \hline
T10 & Đăng xuất rồi đăng nhập bằng tài khoản khác & Dữ liệu thông báo và tác vụ cục bộ không bị lẫn giữa hai tài khoản & Cao \\ \hline
T11 & Lưu body scan và mở lại màn hình hồ sơ hoặc theo dõi & Các chỉ số thể trạng gần nhất như body fat hoặc posture score được cập nhật & Trung bình \\ \hline
\end{tabular}
\normalsize
\end{table}

\subsection{Kiểm thử thông báo và đồng bộ dữ liệu}

Đối với thông báo đẩy, quy trình kiểm thử tập trung vào ba điểm: gửi đúng nội dung tới thiết bị, quan sát hành vi khi ứng dụng ở nền trước hoặc nền sau và kiểm tra việc lưu lại bản ghi thông báo ở phía cục bộ \cite{firebase_fcm}. Đối với phân hệ bài tập, kiểm thử tập trung vào các trạng thái dữ liệu hợp lệ, dữ liệu rỗng, dữ liệu sai chuẩn, thiết bị ngoại tuyến nhưng đã có bộ nhớ đệm và trường hợp tải media không thành công.

\section{Đánh giá kết quả triển khai}

\subsection{Các kết quả đã đạt được}

Qua việc triển khai hiện tại, có thể nhận thấy hệ thống đã hiện thực được hầu hết các thành phần chính của một ứng dụng hỗ trợ tập luyện cá nhân hóa:

\begin{itemize}
    \item có đầy đủ luồng truy cập, khởi tạo hồ sơ và xem trước kế hoạch;
    \item có cấu trúc kế hoạch tập, buổi tập theo lịch và phiên tập tương đối hoàn chỉnh;
    \item có phân hệ tra cứu bài tập theo mô hình ngoại tuyến ưu tiên;
    \item có phân hệ theo dõi tiến độ, hồ sơ, thông báo cục bộ và thông báo đẩy;
    \item có khả năng ghi nhận dinh dưỡng, body scan và AI coach ở mức hỗ trợ trải nghiệm;
    \item có hệ thống nội dung động trên Firebase cùng công cụ hỗ trợ nạp và kiểm tra dữ liệu.
\end{itemize}

\subsection{Đối chiếu yêu cầu và kết quả hiện thực}

Bảng~\ref{tab:requirement-mapping} đối chiếu trực tiếp các nhóm yêu cầu chức năng chính với trạng thái hiện thực trong hệ thống.

\small
\begin{longtable}{|p{1.1cm}|p{3.8cm}|p{2.2cm}|p{5.1cm}|}
\caption{Đối chiếu giữa yêu cầu chức năng và kết quả hiện thực}
\label{tab:requirement-mapping}\\
\hline
\textbf{Mã} & \textbf{Yêu cầu} & \textbf{Trạng thái} & \textbf{Ghi chú} \\ \hline
\endfirsthead
\hline
\textbf{Mã} & \textbf{Yêu cầu} & \textbf{Trạng thái} & \textbf{Ghi chú} \\ \hline
\endhead
\hline
\multicolumn{4}{r}{\textit{Còn tiếp ở trang sau}} \\
\endfoot
\hline
\endlastfoot
F1 & Quản lý truy cập & Đạt & Đã có email/password, Google, chế độ khách, quên mật khẩu và đăng xuất. \\ \hline
F2 & Xác định trạng thái khởi động & Đạt & Màn hình khởi động điều hướng theo trạng thái đăng nhập và khởi tạo hồ sơ. \\ \hline
F3 & Thu thập hồ sơ ban đầu & Đạt & Phần khởi tạo hồ sơ lưu thông tin cá nhân, thiết lập nhắc việc và dữ liệu sở thích tập luyện. \\ \hline
F4 & Sinh kế hoạch tập khởi đầu & Đạt & Có phần xem trước kế hoạch, kế hoạch khởi đầu, lịch buổi tập và nội dung xem trước. \\ \hline
F5 & Quản lý nội dung bài tập & Đạt & Có nhóm bài tập, danh sách bài tập, chi tiết bài tập, trình phát video và đồng bộ dữ liệu mô tả. \\ \hline
F6 & Thực hiện buổi tập & Đạt một phần mạnh & Luồng ghi nhận phiên tập tương đối hoàn chỉnh; buổi tập nhanh còn một số quy tắc cục bộ cần tiếp tục hoàn thiện. \\ \hline
F7 & Theo dõi tiến độ & Đạt một phần mạnh & Có màn hình theo dõi, tổng hợp theo ngày, nhật ký bữa ăn, số đo cơ thể và body scan; một số chỉ số vẫn còn ước lượng. \\ \hline
F8 & Quản lý hồ sơ & Đạt & Có cập nhật hồ sơ, ảnh đại diện, hiển thị thông tin và thiết lập cơ bản. \\ \hline
F9 & Thông báo & Đạt & Có nhắc việc cục bộ, đăng ký nhận thông báo đẩy và lưu vết thông báo trên thiết bị. \\ \hline
F10 & Đồng bộ nội dung động & Đạt & Có kho dữ liệu nội dung động, mô tả cấu trúc dữ liệu và công cụ hỗ trợ nạp dữ liệu ban đầu. \\ \hline
F11 & Hỗ trợ phân tích dinh dưỡng và thể trạng & Đạt một phần mạnh & Đã có meal scan, meal log, body scan và AI hỗ trợ; chưa phải hệ thống chẩn đoán hay tư vấn chuyên sâu. \\ \hline
\end{longtable}
\normalsize

\subsection{Ưu điểm và hạn chế}

Giải pháp hiện tại có các ưu điểm chính là kiến trúc phân lớp rõ ràng, khả năng cập nhật nội dung động và cách tổ chức dữ liệu cục bộ giúp tăng tính ổn định cho trải nghiệm sử dụng. Bên cạnh đó, hệ thống vẫn còn một số điểm cần cải thiện như một số chỉ số theo dõi còn mang tính ước lượng, phân tích dinh dưỡng và thể trạng mới ở mức hỗ trợ, buổi tập nhanh còn phụ thuộc vào một số quy tắc cục bộ và dữ liệu nội dung vẫn cần tiếp tục làm giàu để tăng tính hoàn thiện.

\section{Khả năng triển khai và mở rộng}

Với kiến trúc hiện tại, hệ thống có thể tiếp tục mở rộng theo các hướng như bổ sung thêm nguồn nội dung bài tập, hoàn thiện sâu hơn phần theo dõi sức khỏe, tăng độ phong phú của hệ thống gợi ý, phát triển meal planning theo mục tiêu tăng cơ hoặc giảm mỡ và nâng chất lượng dữ liệu động trên Firebase. Việc tách rõ giao diện, nghiệp vụ và dữ liệu cũng tạo điều kiện thuận lợi hơn cho bảo trì và nâng cấp sau này.

\section{Tổng kết chương}

Chương này đã trình bày ngắn gọn phần cài đặt, kiểm thử và đánh giá hệ thống Fitty. Kết quả hiện tại cho thấy ứng dụng đã hiện thực được các luồng cốt lõi của bài toán, đồng thời vẫn còn không gian để tiếp tục hoàn thiện về dữ liệu, kiểm thử và chiều sâu chức năng trong các giai đoạn phát triển tiếp theo.
